\section{Shared mobility systems}
\label{sec:shared_mobility_systems}

Shared mobility services have become an increasingly prominent component of urban transport systems, offering flexible and cost-efficient alternatives to private vehicle ownership. Broadly defined, shared mobility refers to the short-term, on-demand use of transportation modes that are shared among users, either simultaneously or sequentially~\cite{Shaheen2015}. By optimizing vehicle usage and reducing idle resources, shared mobility combines the flexibility and user autonomy of individualized travel with the efficiency gains of shared infrastructure and logistics. Operating at the intersection of private and public transport, it is increasingly recognized as a key component in the transition toward more sustainable and resilient transport systems~\cite{UNFCCC_2022}, as it reduces reliance on private vehicles and promotes multimodal travel behavior.

Within the landscape of shared mobility, two main operational paradigms can be distinguished based on how the resources are shared~\cite{Shaheen2015, Santi2017}. \textbf{Trip-sharing} aggregates users who share the same \ac{od} pair or route. Examples include carpooling and ride-sharing services. This strategy improves cost efficiency by splitting the cost of a trip among multiple passengers. However, it comes at the expense of flexibility: users must adapt to fixed routes, timetables, or detours. Moreover, the vehicle's capacity can become a limiting factor, particularly during peak demand periods. By contrast, \textbf{vehicle-sharing} involves individual trips using shared vehicles, with no user aggregation. Examples include car-sharing and shared micromobility like scooters, mopeds, and \acp{bss}. These services maximize fleet utilization through trip chaining and offer greater flexibility to users. Nonetheless, they introduce operational challenges, such as vehicle redistribution and availability.

Among these two paradigms, the focus of this thesis lies on vehicle-sharing systems. These are typically characterized by short-term rentals, mobile access via apps or smartcards, and pay-per-use or subscription-based pricing models. Their implementation can follow different deployment models~\cite{Fu2025, NatureEVSharing2023}. In \textbf{station-based systems}, users pick up and return vehicles at designated stations. On the other hand, \textbf{free-floating systems} allow users to locate, unlock, and park vehicles anywhere within a predefined operational area, offering greater spatial flexibility but requiring more sophisticated fleet monitoring and redistribution strategies. Finally, \textbf{hybrid models} combine elements of both, providing station-based hubs while also allowing flexible drop-off within specific zones.

Despite the diversity of shared mobility services, most systems rely on digital infrastructure for real-time tracking, reservations, billing, and operational optimization. These technologies allow providers to monitor demand patterns, adapt fleet deployment, and maintain service quality in dynamic urban contexts. When implemented effectively, vehicle-sharing systems offer several well-documented benefits. For example, station-based car-sharing has been associated with the removal of up to nine private vehicles for each shared car introduced~\cite{Kolleck2021}, while per capita vehicle use has been shown to decrease by 40–60\% among users of such services~\cite{Litman2000}.

Among the different forms of vehicle-sharing, \acp{bss} have become one of the most widespread and emblematic solutions in urban mobility. Their definition, evolution, impacts, and other aspects are explored in the following Sections.